\begin{proofbox}
\par\noindent\begin{minipage}{\linewidth}
  \textbf{\textcolor{CProof}{证明一：门槛最低——只有勾股定理、平方差（初中知识）}}

  \textbf{目标：}不使用正弦、余弦、正切，也不使用任何超出初中知识的定理。
  设 $BC=a$、$CA=b$、$AB=c$。由于三边在公式里的地位相同，可以把\textbf{最长的边}记为 $c$。
  从 $C$ 向 $AB$ 作垂线，垂足为 $D$，令 $CD=h$、$AD=x$，则 $DB=c-x$。
  \textbf{为什么 $D$ 在边 $AB$ 上？} 因为 $AB$ 最长：如果垂足落在 $A$ 的左边，直角三角形 $BCD$ 的斜边 $BC$ 就比 $BD$ 还长，而 $BD>BA$，于是 $BC>AB$，与最长边为 $AB$ 矛盾；落在 $B$ 的右边同理。

  % True 13-14-15 triangle, longest side AB=c.
% T037: exact side-length model, A=(0,0), B=(c,0).
% Arguments a=BC,b=CA,c=AB,scale,mode,BC label,CA label,AB label,caption.
% Modes 0 area, 1 common altitude, 2 tangent pairs, 3 angle, 4 altitude example, 5 area pieces.
\providecommand{\THeronFigure}[9]{%
\begin{center}
\begin{tikzpicture}[scale=#4,font=\small,line join=round,line cap=round]
\pgfmathsetmacro{\tq}{((#2)^2+(#3)^2-(#1)^2)/(2*(#3))}
\pgfmathsetmacro{\th}{sqrt((#2)^2-\tq^2)}
\pgfmathsetmacro{\tp}{((#1)+(#2)+(#3))/2}
\pgfmathsetmacro{\tr}{(#3)*\th/(2*\tp)}
\pgfmathsetmacro{\tix}{((#2)*(#3)+(#3)*\tq)/(2*\tp)}
\coordinate (A) at (0,0);\coordinate (B) at (#3,0);
\coordinate (C) at (\tq,\th);\coordinate (H) at (\tq,0);
\coordinate (I) at (\tix,\tr);
\coordinate (D) at ($(B)!(I)!(C)$);
\coordinate (E) at ($(C)!(I)!(A)$);
\coordinate (F) at ($(A)!(I)!(B)$);
\ifnum#5=0
 \fill[blue!10] (A)--(B)--(C)--cycle;
 \node at ({(#3+\tq)/3},{\th/3}) {$S$};
\fi
\ifnum#5=7
 \fill[blue!10] (A)--(B)--(C)--cycle;
 \pic[draw,angle radius=3mm] {right angle=A--C--B};
\fi
\ifnum#5=1
 \fill[blue!13] (A)--(H)--(C)--cycle;
 \fill[orange!18] (B)--(H)--(C)--cycle;
\fi
\ifnum#5=5
 \fill[blue!13] (I)--(B)--(C)--cycle;
 \fill[orange!18] (I)--(C)--(A)--cycle;
 \fill[green!15!white] (I)--(A)--(B)--cycle;
\fi
\draw[thick] (A)--(B)--(C)--cycle;
\ifnum#5=2\else
 \path (B)--node[pos=.82,right=5pt] {$#6$} (C);
 \path (C)--node[pos=.25,left=5pt] {$#7$} (A);
 \ifnum#5=1\else
  \ifnum#5=4
   \path (A)--node[midway,below=28pt] {$#8$} (B);
  \else
   \ifnum#5=5
    \path (A)--node[midway,below=28pt] {$#8$} (B);
   \else
    \path (A)--node[midway,below=12pt] {$#8$} (B);
   \fi
  \fi
 \fi
\fi
\foreach \v/\where in {A/below left,B/below right,C/above}
 {\fill (\v) circle[radius=1.4pt];\node[\where=4pt] at (\v) {$\v$};}
\ifnum#5=1
 \draw[very thick,CProof] (C)--node[midway,right=4pt] {$h$} (H);
 \pic[draw,angle radius=2.8mm] {right angle=C--H--B};
 \fill (H) circle[radius=1.2pt];\node[below=4pt] at (H) {$D$};
 \path (A)--node[midway,below=16pt] {$x$} (H);
 \path (H)--node[midway,below=16pt] {$c-x$} (B);
\fi
\ifnum#5=4
 \draw[very thick,CProof] (C)--node[midway,right=4pt] {$h$} (H);
 \pic[draw,angle radius=2.8mm] {right angle=C--H--B};
 \fill (H) circle[radius=1.2pt];\node[below=4pt] at (H) {$D$};
\fi
\ifnum#5=3
 \coordinate (K) at ($(B)!(A)!(C)$);
 \draw[CProof,dashed,thick] (A)--node[pos=.55,above=6pt] {$h_a$} (K);
 \pic[draw,angle radius=2.8mm] {right angle=A--K--C};
 \pic[draw,CProof,angle radius=9mm] {angle=A--C--B};
\fi
\ifnum#5=2\THeronIncircleDetails{2}\fi
\ifnum#5=5\THeronIncircleDetails{5}\fi
\ifnum#5=0
 \node[anchor=north,align=center,text width=11cm] at (#3/2,{-1.25/#4}) {#9};
\else
 \node[anchor=north,align=center,text width=11cm] at (#3/2,{-1.65/#4}) {#9};
\fi
\end{tikzpicture}
\end{center}}

\providecommand{\THeronIncircleDetails}[1]{%
\draw[CStatement,thick] (I) circle[radius=\tr];
\draw[gray,dashed] (I)--(A) (I)--(B) (I)--(C);
\draw[very thick,CStatement] (I)--(D) (I)--(E) (I)--node[midway,right=3pt] {$r$} (F);
\pic[draw,angle radius=2mm] {right angle=I--D--C};
\pic[draw,angle radius=2mm] {right angle=I--E--A};
\pic[draw,angle radius=2mm] {right angle=I--F--B};
\foreach \v/\where in {D/above right,E/above left,F/below}
 {\fill (\v) circle[radius=1.3pt];\node[\where=4pt] at (\v) {$\v$};}
\fill (I) circle[radius=1.6pt];\node at ($(I)+(1.3,0)$) {$I$};
\ifnum#1=2
 \draw[very thick,CStatement] (A)--node[midway,left=6pt] {$x$} (E)
  (A)--node[midway,below=16pt] {$x$} (F);
 \draw[very thick,CExample] (B)--node[midway,right=6pt] {$y$} (D)
  (B)--node[midway,below=16pt] {$y$} (F);
 \draw[very thick,CProof] (C)--node[midway,left=6pt] {$z$} (E)
  (C)--node[midway,right=6pt] {$z$} (D);
\fi}

\THeronFigure{13}{14}{15}{0.48}{1}{a}{b}{c}{底边 $AB=c$；两个直角三角形共用同一条高 $CD=h$。}

\end{minipage}\par

  \begin{description}[style=nextline,leftmargin=2.8em,labelsep=0.5em,itemsep=0.55em]
    \item[\textbf{第 1 步：对左右两个直角三角形分别用勾股定理。}]
    \[
      \begin{cases}
       b^2=x^2+h^2, &\text{在 }\triangle ACD\text{ 中},\\
       a^2=(c-x)^2+h^2, &\text{在 }\triangle BCD\text{ 中}.
      \end{cases}
    \]
    \textbf{为什么这样做？}两式都有相同的 $h^2$，相减就可以把不知道的高消掉。

    \item[\textbf{第 2 步：两式相减，先求出 $x$。}]
    \[
      \begin{aligned}
        b^2-a^2
        &=x^2-(c-x)^2\\
        &=x^2-(c^2-2cx+x^2)\\
        &=2cx-c^2.
      \end{aligned}
    \]
    所以
    \[
      2cx=b^2+c^2-a^2,
      \qquad x=\frac{b^2+c^2-a^2}{2c}.
    \]
    \textbf{关键：}没有求出高，却已经把底边的分段 $x$ 表示成了三边。

    \item[\textbf{第 3 步：面积用底乘高，先算面积的平方。}]
    因为 $S=\tfrac12 ch$，所以 $4S^2=c^2h^2$。
    又因为 $h^2=b^2-x^2$，故
    \[
      \begin{aligned}
        16S^2&=4c^2h^2\\
        &=4c^2(b^2-x^2)\\
        &=4b^2c^2-(2cx)^2\\
        &=4b^2c^2-(b^2+c^2-a^2)^2.
      \end{aligned}
    \]
    \textbf{为什么算 $S^2$ 而不是 $S$？}这样可以不急着开根号，先把式子变成两个平方的差。

    \item[\textbf{第 4 步：先用一次平方差，再连续用平方差。}]
    \[
      \begin{aligned}
       16S^2
       &=(2bc)^2-(b^2+c^2-a^2)^2\\
       &=[\,2bc+b^2+c^2-a^2\,]
         [\,2bc-b^2-c^2+a^2\,]\\
       &=[\,(b+c)^2-a^2\,]
         [\,a^2-(b-c)^2\,]\\
       &=(b+c+a)(b+c-a)\\[-1mm]
       &\qquad\cdot(a+b-c)(a-b+c).
      \end{aligned}
    \]
    \textbf{为什么突然出现四个括号？}因为
    $(b+c)^2-a^2=(b+c-a)(b+c+a)$，
    $a^2-(b-c)^2=(a-b+c)(a+b-c)$。
    每次都只是初中学过的 ``$u^2-v^2=(u-v)(u+v)$''。

    \item[\textbf{第 5 步：将四个括号改写成半周长。}]
    令 $p=(a+b+c)/2$，逐个对应：
    \[
      \begin{array}{ll}
       a+b+c=2p,&b+c-a=2(p-a),\\
       a+b-c=2(p-c),&a-b+c=2(p-b).
      \end{array}
    \]
    于是
    \[
      \begin{aligned}
       16S^2
       &=2p\cdot2(p-a)\cdot2(p-b)\cdot2(p-c)\\
       &=16p(p-a)(p-b)(p-c).
      \end{aligned}
    \]
    两边除以 $16$，得到 $S^2=p(p-a)(p-b)(p-c)$。
    因为三角形面积 $S>0$，所以
    \[
      \boxed{S=\sqrt{p(p-a)(p-b)(p-c)}}.
    \]
  \end{description}

  \medskip
\par\noindent\begin{minipage}{\linewidth}
  \textbf{\textcolor{CProof}{证明二：关键步骤最少——内切圆＋两个半角（高中三角函数）}}

  \textbf{目标：}用三个关键环节直接推出 $S^2=p(p-a)(p-b)(p-c)$，看清半周长三个差从哪里来。
  设内切圆圆心为 $I$、半径为 $r>0$，与 $BC$、$CA$、$AB$ 分别切于 $D$、$E$、$F$。
  记 $AF=AE=x$、$BF=BD=y$、$CD=CE=z$。这些相等是因为，从同一个顶点向圆作的两条切线段等长；也可用勾股定理证明，如 $AF^2=AI^2-r^2=AE^2$。

  % D on BC, E on CA, F on AB, matching the formal proof.
% T037: exact side-length model, A=(0,0), B=(c,0).
% Arguments a=BC,b=CA,c=AB,scale,mode,BC label,CA label,AB label,caption.
% Modes 0 area, 1 common altitude, 2 tangent pairs, 3 angle, 4 altitude example, 5 area pieces.
\providecommand{\THeronFigure}[9]{%
\begin{center}
\begin{tikzpicture}[scale=#4,font=\small,line join=round,line cap=round]
\pgfmathsetmacro{\tq}{((#2)^2+(#3)^2-(#1)^2)/(2*(#3))}
\pgfmathsetmacro{\th}{sqrt((#2)^2-\tq^2)}
\pgfmathsetmacro{\tp}{((#1)+(#2)+(#3))/2}
\pgfmathsetmacro{\tr}{(#3)*\th/(2*\tp)}
\pgfmathsetmacro{\tix}{((#2)*(#3)+(#3)*\tq)/(2*\tp)}
\coordinate (A) at (0,0);\coordinate (B) at (#3,0);
\coordinate (C) at (\tq,\th);\coordinate (H) at (\tq,0);
\coordinate (I) at (\tix,\tr);
\coordinate (D) at ($(B)!(I)!(C)$);
\coordinate (E) at ($(C)!(I)!(A)$);
\coordinate (F) at ($(A)!(I)!(B)$);
\ifnum#5=0
 \fill[blue!10] (A)--(B)--(C)--cycle;
 \node at ({(#3+\tq)/3},{\th/3}) {$S$};
\fi
\ifnum#5=7
 \fill[blue!10] (A)--(B)--(C)--cycle;
 \pic[draw,angle radius=3mm] {right angle=A--C--B};
\fi
\ifnum#5=1
 \fill[blue!13] (A)--(H)--(C)--cycle;
 \fill[orange!18] (B)--(H)--(C)--cycle;
\fi
\ifnum#5=5
 \fill[blue!13] (I)--(B)--(C)--cycle;
 \fill[orange!18] (I)--(C)--(A)--cycle;
 \fill[green!15!white] (I)--(A)--(B)--cycle;
\fi
\draw[thick] (A)--(B)--(C)--cycle;
\ifnum#5=2\else
 \path (B)--node[pos=.82,right=5pt] {$#6$} (C);
 \path (C)--node[pos=.25,left=5pt] {$#7$} (A);
 \ifnum#5=1\else
  \ifnum#5=4
   \path (A)--node[midway,below=28pt] {$#8$} (B);
  \else
   \ifnum#5=5
    \path (A)--node[midway,below=28pt] {$#8$} (B);
   \else
    \path (A)--node[midway,below=12pt] {$#8$} (B);
   \fi
  \fi
 \fi
\fi
\foreach \v/\where in {A/below left,B/below right,C/above}
 {\fill (\v) circle[radius=1.4pt];\node[\where=4pt] at (\v) {$\v$};}
\ifnum#5=1
 \draw[very thick,CProof] (C)--node[midway,right=4pt] {$h$} (H);
 \pic[draw,angle radius=2.8mm] {right angle=C--H--B};
 \fill (H) circle[radius=1.2pt];\node[below=4pt] at (H) {$D$};
 \path (A)--node[midway,below=16pt] {$x$} (H);
 \path (H)--node[midway,below=16pt] {$c-x$} (B);
\fi
\ifnum#5=4
 \draw[very thick,CProof] (C)--node[midway,right=4pt] {$h$} (H);
 \pic[draw,angle radius=2.8mm] {right angle=C--H--B};
 \fill (H) circle[radius=1.2pt];\node[below=4pt] at (H) {$D$};
\fi
\ifnum#5=3
 \coordinate (K) at ($(B)!(A)!(C)$);
 \draw[CProof,dashed,thick] (A)--node[pos=.55,above=6pt] {$h_a$} (K);
 \pic[draw,angle radius=2.8mm] {right angle=A--K--C};
 \pic[draw,CProof,angle radius=9mm] {angle=A--C--B};
\fi
\ifnum#5=2\THeronIncircleDetails{2}\fi
\ifnum#5=5\THeronIncircleDetails{5}\fi
\ifnum#5=0
 \node[anchor=north,align=center,text width=11cm] at (#3/2,{-1.25/#4}) {#9};
\else
 \node[anchor=north,align=center,text width=11cm] at (#3/2,{-1.65/#4}) {#9};
\fi
\end{tikzpicture}
\end{center}}

\providecommand{\THeronIncircleDetails}[1]{%
\draw[CStatement,thick] (I) circle[radius=\tr];
\draw[gray,dashed] (I)--(A) (I)--(B) (I)--(C);
\draw[very thick,CStatement] (I)--(D) (I)--(E) (I)--node[midway,right=3pt] {$r$} (F);
\pic[draw,angle radius=2mm] {right angle=I--D--C};
\pic[draw,angle radius=2mm] {right angle=I--E--A};
\pic[draw,angle radius=2mm] {right angle=I--F--B};
\foreach \v/\where in {D/above right,E/above left,F/below}
 {\fill (\v) circle[radius=1.3pt];\node[\where=4pt] at (\v) {$\v$};}
\fill (I) circle[radius=1.6pt];\node at ($(I)+(1.3,0)$) {$I$};
\ifnum#1=2
 \draw[very thick,CStatement] (A)--node[midway,left=6pt] {$x$} (E)
  (A)--node[midway,below=16pt] {$x$} (F);
 \draw[very thick,CExample] (B)--node[midway,right=6pt] {$y$} (D)
  (B)--node[midway,below=16pt] {$y$} (F);
 \draw[very thick,CProof] (C)--node[midway,left=6pt] {$z$} (E)
  (C)--node[midway,right=6pt] {$z$} (D);
\fi}

\THeronFigure{13}{14}{15}{0.48}{2}{a}{b}{c}{同一顶点的两条切线段等长：$AF=AE=x$，$BF=BD=y$，$CD=CE=z$。}

\end{minipage}\par
  \begin{description}[style=nextline,leftmargin=2.8em,labelsep=0.5em,itemsep=0.5em]
    \item[\textbf{第 1 步：用三个切线段表示三边。}]
    \[
      a=y+z,\quad b=z+x,\quad c=x+y,\quad p=x+y+z.
    \]
    故
    \[
      x=p-a,\quad y=p-b,\quad z=p-c.
    \]
    连接 $IA$、$IB$、$IC$，三个以 $a,b,c$ 为底、以 $r$ 为高的小三角形恰好拼成原三角形，因此
    \[
      S=\frac{ar}{2}+\frac{br}{2}+\frac{cr}{2}=pr.
    \]

    \item[\textbf{第 2 步：用两角和的正切公式，只求一条等式。}]
\begin{center}
\begin{tikzpicture}[font=\small,line join=round]
\foreach \origin/\leg/\vertex/\foot/\base/\angle in {0/8/A/F/x/A,4/7/B/F/y/B,8/6/C/E/z/C}{
 \begin{scope}[shift={(\origin,0)}]
  \coordinate (V) at (0,0);\coordinate (H) at ({.32*\leg},0);
  \coordinate (J) at ({.32*\leg},1.28);
  \fill[blue!8] (V)--(H)--(J)--cycle;
  \draw[thick] (V)--node[midway,below=6pt] {$\base$} (H)
    --node[midway,right=6pt] {$r$} (J)--cycle;
  \pic[draw,angle radius=2mm] {right angle=V--H--J};
  \pic[draw,CProof,angle radius=7mm] {angle=H--V--J};
  \node[left=4pt] at (V) {$\vertex$};\node[below right=4pt] at (H) {$\foot$};
  \node[above right=4pt] at (J) {$I$};
  \node at (1.25,2.05) {顶点处半角 $\angle/2$};
  \node at (1.25,-1.05) {$\tan(\angle/2)=r/\base$};
 \end{scope}
}
\end{tikzpicture}
\end{center}

    因 $AI$、$BI$、$CI$ 分别平分三个角，在直角三角形 $AIF$、$BIF$、$CIE$ 中，
    \[
      \tan\frac A2=\frac r x,\quad
      \tan\frac B2=\frac r y,\quad
      \tan\frac C2=\frac r z.
    \]
    又由 $A+B+C=180^\circ$，有 $\tfrac A2+\tfrac B2=90^\circ-\tfrac C2$。
    利用高中学过的两角和公式
    \[
     \tan(\alpha+\beta)=\frac{\tan\alpha+\tan\beta}{1-\tan\alpha\tan\beta},
    \]
    因 $0<A/2+B/2<90^\circ$，正切存在，且两角和公式的分母为正；下面的除法都有意义。
    左边先计算成
    \[
      \tan\left(\frac A2+\frac B2\right)
      =\frac{\frac r x+\frac r y}{1-\frac{r^2}{xy}}
      =\frac{r(x+y)}{xy-r^2}.
    \]
    而右边同一个角的正切是
    \[
      \tan\left(90^\circ-\frac C2\right)
       =\frac1{\tan(C/2)}=\frac z r.
    \]
    两个结果相等：
    \[
      \frac{r(x+y)}{xy-r^2}=\frac z r
      \Longrightarrow r^2(x+y)=xyz-zr^2
      \Longrightarrow \boxed{r^2(x+y+z)=xyz}.
    \]

    \item[\textbf{第 3 步：把 $S=pr$ 代回去。}]
    因 $p=x+y+z$，有 $r^2p=xyz$，从而
    \[
      S^2=p^2r^2=p\cdot xyz
       =p(p-a)(p-b)(p-c).
    \]
    面积取正平方根，得到海伦公式。
  \end{description}
  \textbf{适用提醒：}这种证明很适合理解 ``$p-a$ 就是切线段'' 的来源；若尚未学过正切两角和公式，\textbf{不要把它当作初中证明}。

  \medskip
  \textbf{\textcolor{CProof}{补充：如果只比书写行数，下面三行更短。}}

  已经学过正弦面积公式和余弦定理时，可以把\textbf{后面的常规证明}压缩为
  \[
   \begin{aligned}
    16S^2
    &=4a^2b^2-(a^2+b^2-c^2)^2\\
    &=[(a+b)^2-c^2][c^2-(a-b)^2]\\
    &=(a+b+c)(a+b-c)(a-b+c)(-a+b+c)\\
    &=16p(p-a)(p-b)(p-c).
   \end{aligned}
  \]
  这里第一行来自 ``$S=\tfrac12 ab\sin C$ 加余弦定理''，第二行是\textbf{平方差}；\textbf{它不是另一种新定理}，而是把证明三缩写成考场上的少量步骤。若不清楚第一行为什么成立，请紧接着看完整的常规证明。

  \medskip
\par\noindent\begin{minipage}{\linewidth}
  \textbf{\textcolor{CProof}{证明三：常规证明——正弦面积公式＋余弦定理（高中方法）}}

\THeronFigure{13}{14}{15}{0.48}{3}{a}{b}{c}{两边 $a,b$ 的夹角在顶点 $C$；对应高 $h_a=b\sin C$。}
\end{minipage}\par
  \textbf{第 1 步：先把面积写成两边与夹角。}
  由于 $a=BC$ 与 $b=CA$ 的夹角为 $C$，故
  \[
    S=\frac12 ab\sin C,\qquad S^2=\frac14a^2b^2\sin^2 C.
  \]
  \textbf{第 2 步：余弦定理消去角 $C$。}
  \[
   c^2=a^2+b^2-2ab\cos C
   \quad\Longrightarrow\quad
   \cos C=\frac{a^2+b^2-c^2}{2ab}.
  \]
  由 $\sin^2 C+\cos^2 C=1$，
  \[
   \sin^2C=1-\left(\frac{a^2+b^2-c^2}{2ab}\right)^2.
  \]
  \textbf{第 3 步：代回面积平方，把它拆成两个平方的差。}
  \[
    \begin{aligned}
     16S^2
      &=4a^2b^2\left(1-\frac{(a^2+b^2-c^2)^2}{4a^2b^2}\right)\\
      &=4a^2b^2-(a^2+b^2-c^2)^2\\
      &=(2ab)^2-(a^2+b^2-c^2)^2.
    \end{aligned}
  \]
  \textbf{第 4 步：两次平方差分解，不跳步。}
  \[
    \begin{aligned}
     16S^2
     &=[\,2ab+(a^2+b^2-c^2)\,]
       [\,2ab-(a^2+b^2-c^2)\,]\\
     &=[\,(a+b)^2-c^2\,][\,c^2-(a-b)^2\,]\\
     &=(a+b+c)(a+b-c)(c-a+b)(c+a-b).
    \end{aligned}
  \]
  \textbf{第 5 步：使用 $2p=a+b+c$，把四项分别改写成 $2p$、$2(p-c)$、$2(p-a)$、$2(p-b)$：}
  \[
    16S^2=16p(p-a)(p-b)(p-c).
  \]
  两边除以 $16$，再取正平方根，
  \[
    \boxed{S=\sqrt{p(p-a)(p-b)(p-c)}}.
  \]

  \medskip
\par\noindent\begin{minipage}{\linewidth}
  \textbf{\textcolor{CProof}{延伸结论的证明：为什么内切圆半径 $r=S/p$？}}
\THeronFigure{13}{14}{15}{0.45}{5}{a}{b}{c}{蓝、橙、绿三块分别以 $a,b,c$ 为底，对应高都为 $r$。}
\end{minipage}\par
  连接内心 $I$ 与 $A,B,C$，内心到每条边的垂直距离均为半径 $r$。
  三块小三角形面积分别为 $\tfrac12 ar$、$\tfrac12 br$、$\tfrac12 cr$，所以
  \[
   S=\frac{ar+br+cr}{2}=\frac{a+b+c}{2}r=pr,
   \qquad\boxed{r=\frac Sp}.
  \]
  \textbf{补充：}这只是 ``面积拆成三个同高三角形''，与是否使用海伦公式无关；海伦公式负责\textbf{由三边求 $S$}，$S=pr$ 负责\textbf{由面积与半周长求 $r$}。

  \medskip
  \textbf{\textcolor{CProof}{三种证明怎样选？}}
  \begin{itemize}[leftmargin=*,itemsep=0.25em]
    \item \textbf{想彻底弄懂：}选证明一——只用初中勾股定理和因式分解。
    \item \textbf{想看到半周长差的几何来源：}选证明二——切线段 $x,y,z$ 一眼对应 $p-a,p-b,p-c$。
    \item \textbf{课堂里已学余弦定理：}选证明三——是教材常见的证明路线。
  \end{itemize}
\end{proofbox}
